\documentclass[11pt]{article}
\usepackage[margin=1in]{geometry}
\usepackage{amsmath,amssymb,amsthm}
\usepackage{mathtools}
\usepackage{enumitem}
\usepackage{hyperref}
\usepackage{cleveref}
\usepackage{booktabs}
\usepackage{tabularx}
\usepackage{microtype}
\usepackage{cite}
\usepackage{xurl}
\usepackage{path}
\hypersetup{colorlinks=true,linkcolor=blue,citecolor=blue,urlcolor=blue}
\newtheorem{definition}{Definition}
\newtheorem{lemma}{Lemma}
\newtheorem{remark}{Remark}
\title{Public Request Response Menus for Source-Only Releases\\\large Smallest-Sufficient Follow-Up Handoffs for Basis Tokens, Visible Notices, Canonical Families, Family Choice, Scope, and Completion}
\author{Anonymous}
\date{Unified archive note v0.11 (2026-03-21; archive v474)}
\begin{document}
\maketitle

\begin{abstract}
The source-only branch can now say what a terse paper promises on request, when that promise is complete, whether the paper-level token survives a successor cut, what exact short notice sentence readers saw, what canonical sentence family sits behind that sentence, and why that family was selected for the shipped basis tokens. One operational question still remained implicit: when a referee or reader asks one concrete follow-up about that visible source-only sentence, which already-maintained object is the smallest sufficient thing to hand over? This note gives that paper-facing response choice one home. It defines compact \emph{public request response menus}: notice-keyed maps from common follow-up classes to the smallest sufficient already-owned source-only object.
\end{abstract}

\paragraph{Non-redundancy note.}
Synthesis~56 owns the exact paper-facing source-only contract behind the public sentence, Synthesis~58 owns the family-scoped completion criterion, Synthesis~59 owns the within-release paper-level status token, Synthesis~60 owns the successor carryforward verdict, Synthesis~61 owns the exact outward-facing notice sentence and pointer order, Synthesis~62 owns the canonical sentence family behind that notice, Synthesis~63 owns the token-keyed choice of which such family applies, Synthesis~65 owns the exact predecessor-side bounded response-packet cut that may travel after one menu choice is fixed, Synthesis~66--72 own the successor-maintenance layers beneath that cut, Synthesis~73 owns the request-side terminal stop later terse papers should usually cite, Synthesis~74--75 own the paired terminal / cutover discipline, and Synthesis~17 owns the concrete worked instantiation.\cite{series:publicrequestcontracts,series:requestfulfillmentcertificates,series:publicrequeststatusenvelopes,series:publicrequestcarryforward,series:publicrequestnotices,series:publicrequestnoticenormalforms,series:publicrequestnoticeselections,series:publicrequestresponsepackets,series:publicrequestresponsepacketcaryforward,series:publicrequestresponsepacketdeltaledgers,series:publicrequestresponsepacketrefreshnotices,series:publicrequestresponsepacketrefreshnoticenormalforms,series:publicrequestresponsepacketrefreshnoticeselections,series:publicrequestresponsepacketrefreshresponsemenus,series:publicrequestresponsepacketrefreshresponsepackets,series:publicrequestresponsepacketrefreshresponsepacketclosureverdicts,series:pairedterminalcitationdiscipline,series:pairedterminalcutoverrules,series:workedexample}
This note owns only the follow-up-class selection among those already-maintained paper-facing companions. Later terse papers should cite Synthesis~64 only when the live issue is which already-maintained object is the smallest sufficient handoff for one concrete follow-up about one shipped source-only notice, or which exact response-menu entry selected that handoff. They should stop at Synthesis~59--63 when the live issue is instead the basis token, visible sentence, canonical family, or family-choice owner; widen to Synthesis~65 when they need the exact bounded predecessor packet cut after that handoff is fixed; widen to Synthesis~66--72 when they need successor reuse, packet-local refresh, visible refresh, canonical wrapper, family selection, successor-facing handoff, or the exact successor-facing packet cut; widen to Synthesis~73 when the live issue is whether that request-side chain already closes; widen to Synthesis~74--75 when the remaining issue is paired-terminal reuse or fail-closed cutover; and stop at Synthesis~17 when the concrete maintained worked route is the live issue.

\paragraph{Scope.}
This is not a new public request contract, not a new completion rule, not a new status token, not a new carryforward verdict, not a new visible notice, not a new canonical sentence family, and not a new family-selection rule.
It is a thin follow-up policy note answering: given one shipped source-only notice, what maintained object should be handed over for one concrete follow-up class?

\paragraph{Exports.}
This note exports (i) public request response menus, (ii) a follow-up-class coverage rule, (iii) a shared-follow-up-vocabulary rule, (iv) a smallest-sufficient source-only response rule, (v) a selector-visibility rule saying the immediate owner of each follow-up choice must remain visible beside the selected handoff, (vi) a notice-scoped selection rule, and (vii) an imported-locality rule saying that later terse papers may cite one maintained response menu when they need a stable answer to ``what exactly should I hand over for this kind of follow-up about the source-only notice, and which maintained object selected that handoff?''.

\section{Why response menus deserve their own note}
Once the source-only line owns visible notices, canonical sentence families, and token-keyed family choice, a short paper still faces one last repetition problem.
If a reader asks ``which token makes this sentence true?'', ``which canonical family is this sentence in?'', ``why that family rather than another one?'', ``what exact contract sits behind the on-request promise?'', or ``what completion rule discharges that promise?'', the answer is not always the same artifact.
Without a maintained response menu, later papers either overshare the whole source-only tail or improvise the handoff in prose.
A response menu keeps those follow-up choices explicit and auditable while leaving truth, wording, and family ownership where they already belong.

\begin{definition}[Public request response menu]
Fix one emitted public-request notice $N$ for a source-only paper.
A \emph{public request response menu} is a tuple
\[
M(N)=(\beta,\nu,\phi,\sigma,\kappa,\rho)
\]
containing one imported basis owner $\beta$, one imported emitted-notice owner $\nu$, one imported canonical normal-form owner $\phi$, one imported family-selection owner $\sigma$, one imported public-request-contract owner $\kappa$, and one imported completion-rule owner $\rho$, together with an ordered family of response entries for the follow-up classes
\[
\begin{array}{l}
\textsf{basis-token-check},\ \textsf{visible-sentence-check},\ \textsf{canonical-family-check},\\
\textsf{family-choice-check},\ \textsf{request-scope-check},\ \textsf{completion-rule-check}.
\end{array}
\]
Each entry names the smallest sufficient already-maintained object for that follow-up class together with its canonical citation path and why that object is sufficient.
\end{definition}

\begin{definition}[Follow-up-class coverage rule]
A public request response menu satisfies the \emph{follow-up-class coverage rule} if it covers exactly the six follow-up classes above for the same emitted notice.
So the menu cannot silently omit one of the common paper-facing questions that recur after a source-only sentence is printed.
\end{definition}

\begin{definition}[Shared-follow-up-vocabulary rule]
A public request response menu satisfies the \emph{shared-follow-up-vocabulary rule} if the six repeated follow-up-class labels denote only imported request-handling vocabulary that may recur across many emitted notices, while the selected handoff for each such label remains owned only by the concrete menu for the emitted notice at hand.
So the shared class names do not silently become global owner choices.
\end{definition}

\begin{definition}[Smallest-sufficient source-only response rule]
A public request response menu satisfies the \emph{smallest-sufficient source-only response rule} if it chooses:
\begin{enumerate}[leftmargin=*]
\item the imported status envelope or carryforward envelope for basis-token checks,
\item the imported emitted notice for visible-sentence checks,
\item the imported canonical notice normal form for canonical-family checks,
\item the imported notice selection for family-choice checks,
\item the imported public request contract for request-scope checks, and
\item the imported request-fulfillment-certificate ledger for completion-rule checks.
\end{enumerate}
So the menu changes only follow-up-class selection; it does not widen or replace any imported source-only owner.
\end{definition}

\begin{definition}[Selector-visibility rule]
A public request response menu satisfies the \emph{selector-visibility rule} if every response entry keeps the selecting menu object visible beside the imported selected handoff by recording both the owner-side menu key and the exact selected object key. So a later terse paper can recover not only what should travel, but also which maintained object selected that handoff for the notice in question.
\end{definition}

\begin{definition}[Notice-scoped selection rule]
A public request response menu satisfies the \emph{notice-scoped selection rule} if every response entry is keyed by one emitted notice lineage and the same follow-up-class label may lead to different selected owners only through different concrete menus for different notices.
So the handoff choice stays attached to the imported notice lineage rather than to the shared class label itself.
\end{definition}

\begin{definition}[Imported-locality rule]
A public request response menu satisfies the \emph{imported-locality rule} if every imported key belongs to the same emitted notice lineage and every response entry copies the imported token bundle, rendered sentence, canonical family, family choice, contract cut, or completion entries exactly.
So the menu may not create a new notice sentence, a new contract scope, or a new completion rule while pretending only to select a handoff object.
\end{definition}

\begin{lemma}[Response menus keep source-only follow-up handling brief without changing ownership]
Assume every public request response menu satisfies the follow-up-class coverage rule, the shared-follow-up-vocabulary rule, the smallest-sufficient source-only response rule, the selector-visibility rule, the notice-scoped selection rule, and the imported-locality rule.
Then a later terse paper may cite one menu as the maintained guide for concrete follow-up about one shipped source-only notice without changing the ownership of its basis token, visible sentence, canonical family, family choice, request scope, or completion rule.
\end{lemma}

\begin{proof}
The follow-up-class coverage rule ensures that the common source-only follow-up requests are covered explicitly.
The shared-follow-up-vocabulary rule makes clear that the repeated six class labels are only reused vocabulary rather than hidden global owner choices.
The smallest-sufficient source-only response rule ensures that each request class resolves to one already-owned object and no larger one.
The selector-visibility rule ensures that the menu remains visible as the immediate owner of that choice even after the selected handoff object is reused downstream.
The notice-scoped selection rule keeps that choice attached to the imported notice lineage rather than to the shared class labels themselves.
The imported-locality rule then prevents the menu from rewriting any imported notice, contract, or completion owner.
So the menu changes only paper-facing handoff policy: it stabilizes what to hand over for one follow-up class, but it does not alter the truth or ownership of the underlying source-only line.\qedhere
\end{proof}

\begin{table}[t]
\centering
\small
\begin{tabularx}{\linewidth}{@{}l l X@{}}
\toprule
Follow-up class & Canonical object & Question answered \\
\midrule
Basis-token check & status/carryforward envelope & Which paper-level token or successor-carryforward verdict makes this notice true? \\
Visible-sentence check & emitted notice & Which exact short sentence and first-pointer order did readers actually see? \\
Canonical-family check & notice normal form & Which reusable sentence family and slot bindings does that visible sentence belong to? \\
Family-choice check & notice selection & Why was that canonical family chosen for this imported token combination? \\
Request-scope check & public request contract & Which maintained packet/menu/validator cut sits behind the public ``available on request'' promise? \\
Completion-rule check & request-fulfillment certificates & Which family-scoped completion entries discharge that promised source-only cut? \\
\bottomrule
\end{tabularx}
\caption{A public-request response menu does not add another source-only stage. It chooses the smallest sufficient already-owned object for one concrete follow-up class about a shipped notice.}
\label{tab:publicrequestresponsemenu}
\end{table}

\paragraph{A minimal public response-menu card.}

\begin{table}[t]
\centering
\footnotesize
\begin{tabularx}{\linewidth}{@{}l X@{}}
\toprule
Card field & Worked floor \\
\midrule
Menu keys & \nolinkurl{worked_example_receipt_interlock_within_release_public_request_response_menu} and \nolinkurl{worked_example_receipt_interlock_successor_public_request_response_menu} \\
Basis / emitted-notice floor & Within release: \nolinkurl{worked_example_receipt_interlock_public_request_status_envelope} with emitted notice \nolinkurl{worked_example_receipt_interlock_public_request_status_notice}; successor carryforward: \nolinkurl{worked_example_receipt_interlock_public_request_carryforward_envelope} with emitted notice \nolinkurl{worked_example_receipt_interlock_public_request_carryforward_notice} \\
Shared follow-up vocabulary floor & Exactly six menu classes: \nolinkurl{basis_token_check}, \nolinkurl{visible_sentence_check}, \nolinkurl{canonical_family_check}, \nolinkurl{family_choice_check}, \nolinkurl{request_scope_check}, and \nolinkurl{completion_rule_check} under \nolinkurl{shared_vocabulary_notice_scoped_selection} \\
Canonical owner bindings floor & Both menus bind \nolinkurl{request_scope_check} to \nolinkurl{worked_example_receipt_interlock_public_request_contract} at \nolinkurl{example_public_request_contracts.json} and \nolinkurl{completion_rule_check} to \nolinkurl{worked-example-request-fulfillment-certificates-v1} at \nolinkurl{example_request_fulfillment_certificates.json}; the within-release and successor menus differ only where the visible notice, basis owner, normal form, or selector genuinely differ \\
Selector-visibility floor & Every response entry keeps both the selecting \nolinkurl{public_request_response_menu_key} and the selected \nolinkurl{object_key} visible, so later terse papers can say not only what travels next but which maintained menu chose it \\
Completion-coverage floor & Both menus explicitly cover the same two family certificates: \nolinkurl{public_status_sentence_request_fulfillment_certificate} and \nolinkurl{compact_diff_summary_request_fulfillment_certificate} \\
Escalation pointer & Stop at Synthesis~59--63 when the live issue is the narrower token / wrapper / family / selector owner; widen to Synthesis~65 only when the exact travelled bounded slice is itself live; widen further through Synthesis~66--75 only when successor maintenance or request-side terminal / cutover ownership is already live \\
\bottomrule
\end{tabularx}
\caption{A minimal public response-menu card. This is the smallest public answer later terse papers can quote directly when the live issue is which maintained owner should answer one concrete follow-up about one shipped source-only notice rather than the narrower token/wrapper owners or the wider packet/terminal tail.}
\label{tab:minimal-public-response-menu-card}
\end{table}

This card is intentionally non-decorative: it pins the actual worked menu keys, the exact shared follow-up vocabulary, the concrete request-scope and completion-rule owner bindings, and the selector-visibility rule without silently re-owning the basis token, visible sentence, canonical family, or bounded packet cut that the neighboring notes already own.\cite{series:workedexample}

\paragraph{A forced public response-menu matrix.}
\begin{table}[t]
\centering
\small
\begin{tabularx}{\linewidth}{@{}lX@{}}
\toprule
Field & Forced public answer \\
\midrule
Within-release menu-entry floor & The pair \nolinkurl{worked_example_receipt_interlock_within_release_public_request_response_menu} plus \nolinkurl{request_scope_check} still forces the single selected owner \nolinkurl{worked_example_receipt_interlock_public_request_contract} at \nolinkurl{example_public_request_contracts.json}; the same menu plus \nolinkurl{completion_rule_check} still forces \nolinkurl{worked-example-request-fulfillment-certificates-v1} at \nolinkurl{example_request_fulfillment_certificates.json} \\
Successor-facing menu-entry floor & The pair \nolinkurl{worked_example_receipt_interlock_successor_public_request_response_menu} plus \nolinkurl{request_scope_check} still forces that same contract object, and the same menu plus \nolinkurl{completion_rule_check} still forces that same family-certificate ledger; changing the notice lineage changes the honest menu key, but not the selected owner when the follow-up class is one of these two shared classes \\
Imported-owner floor & Those two forced menu entries remain honest only while the imported basis envelope / emitted notice / canonical family / selector ladder named by the menu stays the same worked ladder for that notice lineage; the menu does not get to silently keep the same answer after any upstream owner has changed \\
Selector-visibility floor & Each forced entry still names both the selecting \nolinkurl{public_request_response_menu_key} and the selected \nolinkurl{object_key}, so later terse papers can recover the exact menu-entry reason rather than citing the contract or certificate ledger as if it chose itself \\
Staleness trigger floor & This matrix goes stale as soon as the emitted notice lineage changes, a different follow-up class becomes live, the selected \nolinkurl{object_key} no longer matches the entry, or the menu ceases to expose the same imported owner ladder beside that entry \\
Escalation pointer & Stop at Synthesis~59--63 when the live issue is still the token / wrapper / family / selector owner; stop here only while one notice-class pair still forces the same selected handoff; widen to Synthesis~65 only after the handoff object is fixed and the live issue becomes the exact bounded travelled slice \\
\bottomrule
\end{tabularx}
\caption{A forced public response-menu matrix. This is the smallest public answer later terse papers can quote directly when the live issue is whether one shipped notice together with one follow-up class still forces the same maintained handoff entry, rather than merely which menu exists in the abstract.}
\label{tab:forced-public-response-menu-matrix}
\end{table}

This matrix is intentionally non-decorative: it pins the actual worked notice-class pairs that still force the contract and completion-ledger handoffs, together with the precise conditions under which that forcing remains honest or goes stale, without silently re-owning the basis token, visible sentence, canonical family, or bounded packet cut that neighboring notes already own.\cite{series:workedexample}

\paragraph{One tiny response-menu drill.}
For the worked successor-facing notice, the pair \nolinkurl{worked_example_receipt_interlock_successor_public_request_response_menu} plus \nolinkurl{request_scope_check} still forces the contract object \nolinkurl{worked_example_receipt_interlock_public_request_contract}. If the live question changes to ``which family-scoped entries discharge that promise?'', the honest stop remains the same menu but the class changes to \nolinkurl{completion_rule_check}, so the forced handoff becomes \nolinkurl{worked-example-request-fulfillment-certificates-v1}. If the live question instead stays on \nolinkurl{request_scope_check} but the visible notice lineage or imported owner ladder changes, the matrix has gone stale and a later terse paper must reopen the exact notice/menu owner rather than keep quoting the old handoff entry. Only if a later paper truly needs the exact bounded file slice that should travel with the already-fixed contract handoff does it widen one note later to Synthesis~65.\cite{series:workedexample}

\section{Worked example pointer}
The maintained worked bundle now ships \nolinkurl{example_public_request_response_menus.json}.\cite{series:workedexample}
It contains exactly two response-menu entries:
\nolinkurl{worked_example_receipt_interlock_within_release_public_request_response_menu} and
\nolinkurl{worked_example_receipt_interlock_successor_public_request_response_menu}.
Those two menu keys share the explicit vocabulary mode
\nolinkurl{shared_vocabulary_notice_scoped_selection}, but their selected handoffs remain notice-scoped.
The worked ledger now states the exact six class-to-owner bindings directly:

\begin{center}
\small
\begin{tabularx}{\linewidth}{@{}l X X@{}}
\toprule
Request class & Within-release selected owner & Successor-facing selected owner \\
\midrule
\nolinkurl{basis_token_check} & \nolinkurl{worked_example_receipt_interlock_public_request_status_envelope} & \nolinkurl{worked_example_receipt_interlock_public_request_carryforward_envelope} \\
\nolinkurl{visible_sentence_check} & \nolinkurl{worked_example_receipt_interlock_public_request_status_notice} & \nolinkurl{worked_example_receipt_interlock_public_request_carryforward_notice} \\
\nolinkurl{canonical_family_check} & \nolinkurl{worked_example_receipt_interlock_within_release_complete_public_request_notice_normal_form} & \nolinkurl{worked_example_receipt_interlock_successor_preserved_complete_public_request_notice_normal_form} \\
\nolinkurl{family_choice_check} & \nolinkurl{worked_example_receipt_interlock_within_release_public_request_notice_selection} & \nolinkurl{worked_example_receipt_interlock_successor_public_request_notice_selection} \\
\nolinkurl{request_scope_check} & \nolinkurl{worked_example_receipt_interlock_public_request_contract} & \nolinkurl{worked_example_receipt_interlock_public_request_contract} \\
\nolinkurl{completion_rule_check} & \nolinkurl{worked-example-request-fulfillment-certificates-v1} & \nolinkurl{worked-example-request-fulfillment-certificates-v1} \\
\bottomrule
\end{tabularx}
\end{center}

So the worked ledger now states directly that notice-to-handoff challenges stop first at the exact menu entry for the relevant visible notice and request class. This note owns only that notice-keyed handoff. Under the paired terminal citation discipline and the paired cutover rule, later terse papers should cite the answer-side stop first and widen into the source-only branch only when the live question has genuinely crossed the answer-side archive cut; once it has crossed, the indispensable local fact here is the exact menu-side owner rather than a rewalk of the whole request branch.\cite{series:pairedterminalcitationdiscipline,series:pairedterminalcutoverrules} In the worked successor-facing branch that owner is exact rather than generic: \nolinkurl{example_public_request_notice.json} $\to$ menu \nolinkurl{worked_example_receipt_interlock_successor_public_request_response_menu} $\to$ the \nolinkurl{request_scope_check} owner \nolinkurl{worked_example_receipt_interlock_public_request_contract}. That is the maintained source-only owner for the public ``available on request'' promise itself. Only if a later paper truly needs the travelled bounded slice rather than that maintained handoff does it widen one note later to the packet cut in Synthesis~65.

\begin{thebibliography}{99}\small

\bibitem{series:publicrequestcontracts}
Anonymous.
\newblock Public Request Contracts for Source-Only Releases: Exact Paper-Level Handoffs from Boilerplate Policy to Maintained Packet Cuts.
\newblock Series synthesis note (Synthesis~56), 2026. (This archive.)

\bibitem{series:requestfulfillmentcertificates}
Anonymous.
\newblock Request Fulfillment Certificates for Source-Only Releases: Exact Completion Criteria for Public Request Contracts and Typed Omitted-Support Cuts.
\newblock Series synthesis note (Synthesis~58), 2026. (This archive.)

\bibitem{series:publicrequeststatusenvelopes}
Anonymous.
\newblock Public Request Status Envelopes for Source-Only Releases: Paper-Level Tokens Compressing Family Completion without Re-owning the Promise.
\newblock Series synthesis note (Synthesis~59), 2026. (This archive.)

\bibitem{series:publicrequestcarryforward}
Anonymous.
\newblock Public Request Carryforward for Source-Only Successor Releases: When Paper-Level Request Status Persists, Advances, or Reopens across Release Evolution.
\newblock Series synthesis note (Synthesis~60), 2026. (This archive.)

\bibitem{series:publicrequestnotices}
Anonymous.
\newblock Public Request Notices for Source-Only Releases: Outward-Facing Sentences for Within-Release Status and Cross-Release Carryforward.
\newblock Series synthesis note (Synthesis~61), 2026. (This archive.)

\bibitem{series:publicrequestnoticenormalforms}
Anonymous.
\newblock Public Request Notice Normal Forms for Source-Only Releases: Canonical Sentence Families for Within-Release Status and Cross-Release Carryforward Notices.
\newblock Series synthesis note (Synthesis~62), 2026. (This archive.)

\bibitem{series:publicrequestnoticeselections}
Anonymous.
\newblock Public Request Notice Selections for Source-Only Releases: Choosing Canonical Notice Families from Within-Release Status and Cross-Release Carryforward Tokens.
\newblock Series synthesis note (Synthesis~63), 2026. (This archive.)

\bibitem{series:publicrequestresponsepackets}
Anonymous.
\newblock Public Request Response Packets for Source-Only Releases: Exact Bounded Handoff Slices after Notice-Keyed Follow-Up Choice.
\newblock Series synthesis note (Synthesis~65), 2026. (This archive.)

\bibitem{series:publicrequestresponsepacketcaryforward}
Anonymous.
\newblock Public Request Response Packet Carryforward for Source-Only Successor Releases: Reuse, Refresh, or Reopen after a Fixed Predecessor Packet Cut.
\newblock Series synthesis note (Synthesis~66), 2026. (This archive.)

\bibitem{series:publicrequestresponsepacketdeltaledgers}
Anonymous.
\newblock Public Request Response Packet Delta Ledgers for Source-Only Successor Releases: Exact Packet-Local Refresh Bindings beneath Reused Source-Only Handoffs.
\newblock Series synthesis note (Synthesis~67), 2026. (This archive.)

\bibitem{series:publicrequestresponsepacketrefreshnotices}
Anonymous.
\newblock Public Request Response Packet Refresh Notices for Source-Only Successor Releases: Exact Visible Reissue Sentences after Packet-Local Refresh.
\newblock Series synthesis note (Synthesis~68), 2026. (This archive.)

\bibitem{series:publicrequestresponsepacketrefreshnoticenormalforms}
Anonymous.
\newblock Public Request Response Packet Refresh Notice Normal Forms for Source-Only Successor Releases: Canonical Visible Wrapper Families after Packet-Local Refresh.
\newblock Series synthesis note (Synthesis~69), 2026. (This archive.)

\bibitem{series:publicrequestresponsepacketrefreshnoticeselections}
Anonymous.
\newblock Public Request Response Packet Refresh Notice Selections for Source-Only Successor Releases: Choosing Visible Wrapper Families after Packet-Local Refresh.
\newblock Series synthesis note (Synthesis~70), 2026. (This archive.)

\bibitem{series:publicrequestresponsepacketrefreshresponsemenus}
Anonymous.
\newblock Public Request Response Packet Refresh Response Menus for Source-Only Successor Releases: Smallest-Sufficient Successor Handoffs after Visible Packet Refresh.
\newblock Series synthesis note (Synthesis~71), 2026. (This archive.)

\bibitem{series:publicrequestresponsepacketrefreshresponsepackets}
Anonymous.
\newblock Public Request Response Packet Refresh Response Packets for Source-Only Successor Releases: Exact Successor-Facing Packet Cuts after Menu Selection.
\newblock Series synthesis note (Synthesis~72), 2026. (This archive.)

\bibitem{series:publicrequestresponsepacketrefreshresponsepacketclosureverdicts}
Anonymous.
\newblock Public Request Response Packet Refresh Response Packet Closure Verdicts for Source-Only Successor Releases: Stable Request-Side Terminal Forms after Visible Refresh.
\newblock Series synthesis note (Synthesis~73), 2026. (This archive.)

\bibitem{series:pairedterminalcitationdiscipline}
Anonymous.
\newblock Paired Terminal Citation Discipline for Stable Review Spines: One Answer-Side Stop and One Request-Side Capstone for Terse Papers.
\newblock Series synthesis note (Synthesis~74), 2026. (This archive.)

\bibitem{series:pairedterminalcutoverrules}
Anonymous.
\newblock Paired Terminal Cutover Rules for Stable Review Spines: When to Stop at the Checked Answer, When to Widen to Source-Only, and When to Fail Closed to the Reopening Owner.
\newblock Series synthesis note (Synthesis~75), 2026. (This archive.)

\bibitem{series:workedexample}
Anonymous.
\newblock Worked Example: Receipt Interlock for an Anonymous-DHT Reader-Privacy Claim.
\newblock Series synthesis note (Synthesis~17), 2026. (This archive.)

\end{thebibliography}
\end{document}
